\documentclass[11pt]{article}

\usepackage[T1]{fontenc}
\usepackage[utf8]{inputenc}
\usepackage[margin=25mm]{geometry}
\usepackage{amsmath}
\usepackage{booktabs}
\usepackage{graphicx}
\usepackage{microtype}
\usepackage{siunitx}
\usepackage{placeins}
\usepackage{hyperref}

\hypersetup{colorlinks=true,linkcolor=blue,urlcolor=blue}
\sisetup{detect-all}
\graphicspath{{figures/}}

\title{Camera-Parameter Uncertainty: Zoom5 Visual Overview}
\author{}
\date{}

\begin{document}
\maketitle

\noindent
This note follows the current zoom5 camera-to-droplet analysis. It uses the
saved checkerboard bootstrap, the saved ArUco corner fit, and the virtual
pendant-drop contour with the drop assumed to lie in the marker plane
(\(\delta=0\)). The plotting program
\texttt{ProjectScripts/figures\_camera\_uncertainty.jl} reads those saved
results and writes the PDFs in \texttt{docs/figures/}; it does not calibrate,
bootstrap, or run Monte Carlo.

\noindent\textit{Implementation pointers:} intrinsic bootstrap
\texttt{src/camera\_uncertainties.jl:374}; ArUco corner loading and fitting
\texttt{src/camera\_uncertainties.jl:668}; inverse projection and its Jacobian
\texttt{src/camera\_uncertainties.jl:1003, 1173}; normal covariance
\texttt{src/camera\_uncertainties.jl:1340, 1356}; relative-position analysis
\texttt{src/camera\_uncertainties.jl:1621}.

\section{Intrinsic parameters}

The reference intrinsics come from the zoom5 checkerboard calibration in
\texttt{camint\_all.jld2}. It contains 49 calibration images with 117
detected checkerboard corners per image. The camera model has four pinhole
parameters and five distortion coefficients,
\[
  \boldsymbol{\phi}
  = (f_x,f_y,c_x,c_y,k_1,k_2,p_1,p_2,k_3)^T.
\]
For each observed checkerboard point, the code subtracts the point projected
by the fitted calibration from the observed image coordinate. Pooling the
resulting 5733 two-dimensional residuals gives the residual covariance
\[
  \Sigma_{\mathrm{corner}}
  = \operatorname{cov}
    \begin{pmatrix}e_x\\e_y\end{pmatrix}
  \simeq
  \begin{pmatrix}
    6.646 & -0.361\\
    -0.361 & 3.936
  \end{pmatrix}\ {\rm px}^2.
\]
The off-diagonal entry retains the measured x--y residual correlation; this
is a pooled covariance, not a separate covariance estimated for every
checkerboard corner.

The cached 500-draw bootstrap perturbs each detected checkerboard corner
independently with noise from
\(\mathcal{N}(0,\Sigma_{\mathrm{corner}})\), then recalibrates the camera.
The sample covariance of the resulting parameter vectors is
\[
  \Sigma_{\phi}
  = \frac{1}{499}\sum_{b=1}^{500}
    (\boldsymbol{\phi}^{(b)}-\bar{\boldsymbol{\phi}})
    (\boldsymbol{\phi}^{(b)}-\bar{\boldsymbol{\phi}})^T.
\]
The bootstrap cache stores its residuals, \(\Sigma_{\mathrm{corner}}\),
500-by-9 `samples`, the bootstrap mean vector \(\bar{\phi}\), and covariance
\(\Sigma_\phi\). The bootstrap implementation does
not set an explicit random seed; this overview uses the existing cache and
does not regenerate it. The separate analysis seed controls the later
camera-to-droplet draws, not this cached calibration bootstrap.

Figure~\ref{fig:intrinsic-sigma} reports relative standard deviations for
\(f_x,f_y,c_x,c_y\) and absolute coefficient standard deviations for
\(k_1,\ldots,k_3\). Figures~\ref{fig:intrinsic-correlation} and
\ref{fig:intrinsic-pairs} show parameter dependence: diagonal entries in the
heatmap are one, while the scatter panels expose the shape and sign of the
strongest bootstrap correlations.

\begin{figure}[htbp]
  \centering
  \includegraphics[width=\textwidth]{A1_intrinsic_marginal_sigma.pdf}
  \caption{Marginal intrinsic standard deviations from \(\Sigma_\phi\).
  Focal lengths and principal points are shown as relative standard
  deviations in percent; distortion coefficients are shown as absolute
  standard deviations.}
  \label{fig:intrinsic-sigma}
\end{figure}

\begin{figure}[htbp]
  \centering
  \includegraphics[width=0.88\textwidth]{A2_intrinsic_correlation.pdf}
  \caption{Annotated correlation matrix of the cached 500-draw intrinsic
  bootstrap. Red and blue indicate positive and negative correlation;
  stronger color means larger absolute correlation.}
  \label{fig:intrinsic-correlation}
\end{figure}

\begin{figure}[htbp]
  \centering
  \includegraphics[width=\textwidth]{A3_intrinsic_pair_scatter.pdf}
  \caption{Bootstrap scatter for the four most strongly correlated intrinsic
  pairs: \(f_x\)--\(f_y\), \(k_2\)--\(k_3\), \(k_1\)--\(k_2\), and
  \(k_1\)--\(k_3\). Focal and principal-point axes use relative deviations;
  distortion axes use absolute coefficient deviations.}
  \label{fig:intrinsic-pairs}
\end{figure}

\FloatBarrier
\section{Extrinsic parameters from the ArUco marker}

The analysis uses one zoom5 marker image and the fitted corners saved in
\texttt{aruco\_corner\_results\_zoom5\_camtodrop.jld2}. For each detected
OpenCV corner, \texttt{fit\_corner} samples six intensity profiles along
each of the two adjoining edges, starting 8 pixels from the corner. A
sigmoid is fitted to each profile; its center estimates the local edge
position. The edge positions are then regressed against distance along that
edge. If \(X=[\boldsymbol{1}\ \ s]\) is the line design matrix and \(e\) its
residual vector, the code estimates the residual variance as
\[
  \hat{s}^{\,2} = \frac{e^T e}{m-2},
  \qquad
  \widehat{\operatorname{cov}}(a,b)
  = \hat{s}^{\,2}(X^TX)^{-1}.
\]
Only the fitted intercept variance,
\(\widehat{\operatorname{var}}(a)
=\hat{s}^{\,2}[(X^TX)^{-1}]_{11}\), is propagated for each edge; the fitted
slopes are treated as fixed in the corner-intersection covariance.

For the two edge tangents \(t_1,t_2\), normals \(n_1,n_2\), fitted slopes
\(b_1,b_2\), and edge intercepts \(a_1,a_2\), the implementation forms
\[
  A =
  \begin{pmatrix}
    (n_1-b_1t_1)^T\\
    (n_2-b_2t_2)^T
  \end{pmatrix},
  \qquad
  q=A^{-1}\begin{pmatrix}a_1\\a_2\end{pmatrix},
\]
and returns the corner-correction covariance
\[
  \Sigma_q
  = A^{-1}
    \begin{pmatrix}
      \operatorname{var}(a_1)&0\\
      0&\operatorname{var}(a_2)
    \end{pmatrix}
    A^{-T}.
\]
Thus, each saved \(2\times2\) covariance describes the x/y uncertainty of
one fitted correction relative to its OpenCV corner. The calculation
assumes the two edge-intercept errors are independent; the four corner
covariances are also sampled independently by the later pose bootstrap.

The fixed-intrinsics pose bootstrap holds the nominal camera matrix and
distortion fixed, perturbs each of the four fitted corners using its own
saved \(2\times2\) covariance, and runs solvePnP for 1000 draws. The marker
side length is taken as exactly \(\SI{10}{\milli\metre}\). The pose mean is
summarized on the rotation manifold; its six-dimensional samples are
\[
  \boldsymbol{\eta}_i =
  \left(
    \log(R_{\rm mean}^T R_i),\ t_i-t_{\rm mean}
  \right),
  \qquad
  \Sigma_{\rm ext}=\operatorname{cov}(\boldsymbol{\eta}).
\]
The full pose covariance \(\Sigma_{\rm ext}\) is saved; rotations are right-perturbation
radians and translations are millimetres. The analysis output also stores
the 1000 fixed-intrinsics rotations and translations themselves.

Figure~\ref{fig:aruco-corners} shows the source image, fitted corners, and
the covariance ellipses. The ellipses are enlarged 20 times for visibility;
their orientation and relative axis lengths come from the fitted
covariances, but their displayed size is not the physical one. Figures
\ref{fig:pose-sigma} and \ref{fig:pose-correlation} summarize the fixed
intrinsics pose bootstrap.

\subsection{Corner-noise assumptions}

The table gives the saved per-corner covariance summaries. “Previous” and
“next” refer to the marker-corner order; the angles are the acute angles
between the ellipse major axis and each adjoining, unoriented edge direction.
\begin{table}[htbp]
  \centering
  \small
  \begin{tabular}{c c c c c c}
    \toprule
    Corner & \(\sigma_x\) [px] & \(\sigma_y\) [px] & \(\rho_{xy}\) &
    Major--previous & Major--next \\
    \midrule
    1 & 0.823 & 1.569 & \(+0.038\) & \(0.95^\circ\) & \(89.03^\circ\) \\
    2 & 1.144 & 1.478 & \(-0.017\) & \(87.48^\circ\) & \(2.12^\circ\) \\
    3 & 1.050 & 0.355 & \(-0.008\) & \(89.95^\circ\) & \(0.18^\circ\) \\
    4 & 1.279 & 0.449 & \(-0.015\) & \(0.01^\circ\) & \(89.71^\circ\) \\
    \bottomrule
  \end{tabular}
  \caption{Marginal standard deviations, within-corner x--y correlation, and
  major-axis orientation of the saved 2-by-2 corner-fit covariances.}
  \label{tab:corner-covariance}
\end{table}

A corner is the intersection of two fitted edge lines, so independent errors
in the two line offsets are transformed by the line-intersection geometry
into an x--y covariance; non-axis-aligned edges can therefore produce an
off-diagonal term. Here the marker edges are nearly aligned with the pixel
axes, and all four measured correlations are small, but the full covariance
matrices are retained.

Figure~\ref{fig:corner-noise-pixels} shows 500 reproducible perturbations per
corner, sampled for plotting only with a fixed seed. The 1\(\sigma\) and
2\(\sigma\) ellipses are contours of Mahalanobis radius 1 and 2; every panel
uses the same pixel scale and marks the two incident edge directions.

Both pose-sampling paths use each corner's full \(2\times2\) covariance:
the fixed-intrinsics bootstrap applies the multivariate draw to each
corner, and the joint loop does the same for every paired intrinsic sample.
In particular, the sampled noise is not formed from the diagonal standard
deviations alone.

Both \texttt{Camera\_extrinsics.jl} and
\texttt{uncertainty\_camtodrop.jl} use six equally spaced profile centers per
edge, from 8 pixels to \(0.4\) times the shortest adjacent edge; each profile
extends \(\min(40,\text{shortest edge}/12)\) pixels to either side of the
edge. For the saved image, the ranges from the two endpoints of each edge
are:
\begin{center}
  \small
  \begin{tabular}{c c c c}
    \toprule
    Marker edge & Length [px] & From first endpoint [px] &
    From second endpoint [px] \\
    \midrule
    1--2 & 1284.1 & corner 1: 8--504.0 & corner 2: 8--506.4 \\
    2--3 & 1266.0 & corner 2: 8--506.4 & corner 3: 8--506.4 \\
    3--4 & 1275.0 & corner 3: 8--506.4 & corner 4: 8--504.0 \\
    4--1 & 1260.1 & corner 4: 8--504.0 & corner 1: 8--504.0 \\
    \bottomrule
  \end{tabular}
\end{center}
These windows do not overlap: neighboring endpoint fits leave gaps of about
250--274 pixels along the shared edge. Given independent local pixel-fit
noise, using separate windows supports treating the random fit errors as
independent. It does not rule out common systematic errors across the
corners.

Figure~\ref{fig:corner-correlation} displays the \(8\times8\) model
correlation matrix in \((x_1,y_1,\ldots,x_4,y_4)\) order. Its diagonal
\(2\times2\) blocks are the four fitted within-corner covariances; its
off-diagonal blocks are zero because the current model assumes independence
between corners. Correlations estimated from repeated stationary-marker
images could be compared directly with this matrix: compare the diagonal
blocks with the saved per-corner covariances and the empirical
off-diagonal blocks with zero, after accounting for uncertainty and any
residual whole-image motion.

The current treatment covers random fit noise: two edge-intercept variances
are propagated into each corner covariance, and each corner is sampled
independently. It does not model errors shared by all corners, such as blur,
illumination changes, edge-placement bias, or ink spread. A repeated set of
images of the same stationary marker could estimate cross-corner covariance:
after registering frames and separating any whole-image pose motion, form
the eight-component residual vector of all four corners per image and
estimate its \(8\times8\) covariance. Its off-diagonal \(2\times2\) blocks
would quantify correlated motion between corners. The data folder currently
contains only one zoom5 marker image, \texttt{data/test\_droplets/zoom5.JPG};
there are no repeated stationary-marker images to analyze.

For a deterministic sensitivity check, the pose is solved once from the
unshifted fitted corners and again after applying the indicated shift to all
four corners. The table reports right-rotation changes, camera-frame
translation changes (including \(t_z\)), median absolute normal displacement
of the same virtual contour, and the change in its apex-radius measure. No
Monte Carlo draws are added. The comparison scale is the existing median
joint-camera normal standard deviation, \(6.02\,\mu\mathrm{m}\).
\begin{table}[htbp]
  \centering
  \small
  \begin{tabular}{l c c c c c}
    \toprule
    Case & \(\Delta\boldsymbol{\theta}\) [mrad] &
    \(\Delta\boldsymbol{t}\) [\(\mu\)m] &
    Median \(|\Delta n|\) [\(\mu\)m] &
    \(\Delta R_{\rm apex}\) [\(\mu\)m] &
    Camera \(\sigma\) \\
    \midrule
    a: \(+0.1\) px in x & \((-0.001,-0.002,0.000)\) &
    \((0.79,0.00,0.00)\) & 0.73 & \(-2.70\) & \(0.12\times\) \\
    b: \(+0.1\) px in y & \((-0.001,0.000,0.000)\) &
    \((0.00,0.79,-0.01)\) & 0.31 & \(+0.03\) & \(0.05\times\) \\
    c: \(0.1\) px outward & \((0.002,-0.002,0.000)\) &
    \((-0.86,0.07,-4.51)\) & 0.73 & \(+3.01\) & \(0.12\times\) \\
    \bottomrule
  \end{tabular}
  \caption{Deterministic pose and contour response to shared corner shifts.
  Translation components are \((t_x,t_y,t_z)\); “outward” shifts each corner
  radially from the four-corner centroid, making the marker appear larger.}
  \label{tab:shared-corner-sensitivity}
\end{table}

\begin{figure}[htbp]
  \centering
  \includegraphics[width=0.88\textwidth]{B1_aruco_corner_uncertainty.pdf}
  \caption{Zoom5 ArUco marker crop with the four fitted corners and their
  \(2\times2\) covariance ellipses. Each unit-covariance ellipse is enlarged
  by a factor of 20; this is a display enlargement, not a claim that the
  fitted corner uncertainty is 20 times larger. The lime arrows indicate the
  directions of the fitted marker edges.}
  \label{fig:aruco-corners}
\end{figure}

\begin{figure}[htbp]
  \centering
  \includegraphics[width=0.92\textwidth]{B4_corner_noise_in_pixels.pdf}
  \caption{Corner-fit perturbations in pixel coordinates. Each panel shows
  500 fixed-seed draws from that corner's full \(2\times2\) covariance, the
  1\(\sigma\) and 2\(\sigma\) Mahalanobis ellipses, the two incident marker
  edge directions, and \(\sigma_x,\sigma_y,\rho_{xy}\). All panels use equal
  aspect ratio and identical axis limits. The random seed is used only for
  plotting and does not change saved analysis results.}
  \label{fig:corner-noise-pixels}
\end{figure}

\begin{figure}[htbp]
  \centering
  \includegraphics[width=0.82\textwidth]{B5_corner_correlation_matrix.pdf}
  \caption{Correlation matrix of all eight corner coordinates in the model.
  The outlined diagonal blocks are the within-corner correlations; the
  zero off-diagonal blocks encode the independence assumption between
  corners.}
  \label{fig:corner-correlation}
\end{figure}

\begin{figure}[htbp]
  \centering
  \includegraphics[width=0.78\textwidth]{E1_shared_corner_systematic_sensitivity.pdf}
  \caption{Median absolute normal displacement of the virtual contour from
  the three deterministic shared-corner shifts. The dashed line marks the
  median joint-camera normal \(\sigma\) for comparison; all three shifts are
  below that scale.}
  \label{fig:shared-corner-sensitivity}
\end{figure}

\begin{figure}[htbp]
  \centering
  \includegraphics[width=0.9\textwidth]{B2_pose_marginal_sigma.pdf}
  \caption{Marginal standard deviations from the full fixed-intrinsics pose
  covariance. Right-rotation components are expressed in milliradians and
  translations in micrometres.}
  \label{fig:pose-sigma}
\end{figure}

\begin{figure}[htbp]
  \centering
  \includegraphics[width=0.66\textwidth]{B3_pose_correlation.pdf}
  \caption{Correlation matrix of the six-dimensional fixed-intrinsics pose
  covariance. The separator distinguishes rotational from translational
  components; off-diagonal terms within and across these groups are retained.}
  \label{fig:pose-correlation}
\end{figure}

\FloatBarrier
\section{Propagation assuming independence}

Let \(\boldsymbol{p}=(\boldsymbol{\phi},\boldsymbol{\eta})\) contain the
nine intrinsic parameters and six pose increments. The independent model
sets their cross-covariance to zero:
\[
  \Sigma_{\rm ind} =
  \begin{pmatrix}
    \Sigma_\phi & 0\\
    0 & \Sigma_{\rm ext}
  \end{pmatrix}.
\]
For the contour output \(x=(r_1,z_1,\ldots,r_N,z_N)\), linearization gives
\[
  \Sigma_{x,\rm ind}
  =J_{\rm ind}\Sigma_{\rm ind}J_{\rm ind}^T
  =J_\phi\Sigma_\phi J_\phi^T
   +J_\eta\Sigma_{\rm ext}J_\eta^T.
\]
The image-pixel contribution is tracked separately as
\(\Sigma_{\rm pix}=\sigma_{\rm pix}^2J_{\rm pix}J_{\rm pix}^T\), with the
placeholder \(\sigma_{\rm pix}=0.3\) px. The 500-draw independent Monte
Carlo check samples the same block-diagonal 15-parameter Gaussian and sends
each draw through the nonlinear inverse projection. Its median normal
standard deviation is \(\SI{28.43}{\micro\metre}\), close to the
linearized \(\SI{28.50}{\micro\metre}\).

The right panel of Figure~\ref{fig:joint-correlation} is this independent
parameter model. Its intrinsic--pose cross-block is zero by construction;
the intrinsic--intrinsic and pose--pose blocks still contain their measured
correlations.

\FloatBarrier
\section{Joint propagation}

For the joint calculation, 2000 intrinsic vectors are drawn from the
Gaussian defined by the nominal intrinsics and \(\Sigma_\phi\). Each vector
is paired, by row, with one set of independently perturbed ArUco corners.
solvePnP is rerun using that draw's focal lengths, principal point, and
distortion coefficients. The resulting pose increment is stored as
\[
  \boldsymbol{\eta}_i =
  \left(\log(R_{\rm mean}^T R_i),\ t_i-t_{\rm mean}\right).
\]
The saved `joint_intrinsic_samples` and `joint_pose_increments` therefore
preserve the actual pairing. The sample covariance of their concatenation
provides the left-panel 15-by-15 joint correlation matrix in
Figure~\ref{fig:joint-correlation}. The independent and joint matrices are
shown side by side; black rules separate intrinsic and pose blocks.

Changing the focal length changes the marker's predicted image size, and
solvePnP responds by changing its estimated depth. These changes can move
together so that the ratio of focal length to camera-marker distance, which
controls projected scale, changes less than either parameter alone. In this
run, \(\rho(f_x,\delta t_z)=0.647\) and
\(\rho(f_y,\delta t_z)=0.649\). The strongest intrinsic--pose correlations
are \(\rho(c_y,\delta t_y)=-0.983\) and
\(\rho(c_x,\delta t_x)=-0.976\); they are included in Figure~\ref{fig:joint-pairs}.
These are correlations induced by the specified solvePnP procedure, not
shared checkerboard and marker image noise: the two corner-noise sources
remain independent in this model.

\begin{figure}[htbp]
  \centering
  \includegraphics[width=\textwidth]{C1_C2_joint_vs_independent_correlation.pdf}
  \caption{Side-by-side 15-by-15 parameter correlation matrices. Left:
  paired joint draws. Right: the independent model with zero
  intrinsic--pose cross-correlation. The top-left blocks are intrinsic
  parameters, the bottom-right blocks are right-pose increments, and the
  black separators mark their boundary.}
  \label{fig:joint-correlation}
\end{figure}

\begin{figure}[htbp]
  \centering
  \includegraphics[width=\textwidth]{C3_intrinsic_pose_coupling.pdf}
  \caption{Independent parameter draws (gray) overlaid with joint paired
  draws (orange) for the six strongest intrinsic--pose correlations.
  \(f_x\)--\(\delta t_z\) is included among them
  (\(\rho=0.647\)); principal-point/translation pairs show stronger
  coupling. Focal lengths and principal points are plotted as relative
  deviations, distortion as absolute coefficient deviations, rotations in
  milliradians, and translations in micrometres.}
  \label{fig:joint-pairs}
\end{figure}

\FloatBarrier
\section{Effect on the droplet contour}

The analysis generates a virtual \texttt{PD\_fn} pendant-drop contour with
100 resampled points and projects it through the nominal zoom5 camera and
marker pose. It is not a segmentation of a measured droplet. At
\(\delta=0\), the median normal-direction standard deviations are:
\[
\begin{array}{ll}
\text{independent camera:} & \SI{28.50}{\micro\metre},\\
\text{joint camera:} & \SI{6.02}{\micro\metre},\\
\text{pixel contribution:} & \SI{2.37}{\micro\metre},\\
\text{joint camera plus pixel:} & \SI{6.47}{\micro\metre}.
\end{array}
\]
Thus the joint camera median is about 21\% of the independent camera
median. This is the uncertainty reduction expected when the intrinsic--pose
correlations are retained; it is not obtained by changing the pixel-noise
assumption. Figure~\ref{fig:normal-sigma} shows the pointwise contributions.

For visibility, Figure~\ref{fig:contour-bands} enlarges both normal bands
10 times. Figure~\ref{fig:sample-contours} overlays 30 deterministic
subsets from each saved Monte Carlo ensemble, with deviations from the
nominal contour enlarged 10 times. The plotted subsets are for visual
comparison only; statistics below use all 500 independent or 2000 joint
samples.

The apex-radius measure has mean and standard deviation
\(\SI{1.522}{\milli\metre}\pm\SI{0.100}{\milli\metre}\) for the independent
model and
\(\SI{1.528}{\milli\metre}\pm\SI{0.026}{\milli\metre}\) for the joint model.
Figure~\ref{fig:apex-radius} shows the corresponding full histograms.
Figure~\ref{fig:modes} compares the leading contour-covariance modes,
projected onto rigid translation/rotation, scale, and residual shape
components. Each plotted mode is normalized to 100\%; modes with nearly
equal eigenvalues can rotate within their shared subspace, so their
individual component bars should not be over-interpreted.

\begin{figure}[htbp]
  \centering
  \includegraphics[width=0.92\textwidth]{D1_normal_sigma_components.pdf}
  \caption{Normal-direction standard deviation along normalized contour
  arclength: independent camera, joint camera, pixel contribution, and joint
  camera plus pixel. Values are plotted in micrometres.}
  \label{fig:normal-sigma}
\end{figure}

\begin{figure}[htbp]
  \centering
  \includegraphics[width=0.82\textwidth]{D2_contour_uncertainty_bands.pdf}
  \caption{Virtual contour with independent and joint normal-direction
  one-standard-deviation bands. Both bands are enlarged by a factor of 10
  for visibility; the contour itself is not enlarged.}
  \label{fig:contour-bands}
\end{figure}

\begin{figure}[htbp]
  \centering
  \includegraphics[width=0.82\textwidth]{D3_sampled_contours.pdf}
  \caption{Thirty independent and thirty joint sampled contours around the
  same virtual baseline. Deviations are multiplied by 10 for display. The
  samples are selected deterministically from the saved 500- and 2000-draw
  ensembles.}
  \label{fig:sample-contours}
\end{figure}

\begin{figure}[htbp]
  \centering
  \includegraphics[width=0.82\textwidth]{D4_apex_radius_histogram.pdf}
  \caption{Apex-radius measure distributions from all independent and joint
  contour draws. The radius axis is expressed in micrometres; the two
  ensembles contain 500 and 2000 draws, respectively.}
  \label{fig:apex-radius}
\end{figure}

\begin{figure}[htbp]
  \centering
  \includegraphics[width=\textwidth]{D5_covariance_modes.pdf}
  \caption{Rigid, scale, and residual-shape fractions for the five leading
  camera covariance modes, computed separately for the independent and
  joint camera contour covariances. Bars within each mode show the projected
  fraction of that eigenvector; colors are consistent across panels.}
  \label{fig:modes}
\end{figure}

\FloatBarrier
\section{Assumptions and limitations}

\begin{itemize}
  \item The drop is assumed to lie in the marker plane,
  \(\delta=0\). The separately computed relative-position assessment gives
  \(Z=40.35\) mm and a break-even depth offset of about \(0.268\) mm for this
  mount; a nonzero offset uncertainty is not included in the plotted
  camera-only comparison.
  \item The contour is a virtual \texttt{PD\_fn} contour, not an observed
  boundary. Results therefore compare propagation models for this contour,
  not measurement error against a known physical shape.
  \item The pixel uncertainty \(\sigma_{\rm pix}=0.3\) px is a placeholder,
  not a measured contour-noise level.
  \item The four ArUco corners are perturbed independently using their
  respective \(2\times2\) fitted covariances. Cross-corner correlations are
  not estimated by \texttt{fit\_corner}.
  \item Marker side length is treated as exact, and the pose covariance is
  based on one marker in one zoom5 image. Other markers, images, and zoom
  settings are not represented.
  \item The intrinsic bootstrap and ArUco corner noise are independent.
  The joint treatment captures the dependence created when the sampled
  intrinsics are used in solvePnP; it does not jointly resample checkerboard
  calibration images and marker observations.
  \item The intrinsic bootstrap implementation does not explicitly seed its
  NumPy random generator. This report uses the saved 500-draw cache rather
  than claiming the calibration bootstrap can be recreated from the
  camera-to-droplet seed.
\end{itemize}

\end{document}
